\paragraph{The specified second element at the next local norm test.}
The successful two-step family in A.156 also gives a literal
example of a representative that cannot be used unchanged at the
following step.  Let \(K_3=\Q(\rho_3)\) with
\(\rho_3^3=\rho_2\), in the same fixed tower.
Then the particular \(\beta_2=-\gamma\) in (A.156.4) satisfies
\[
                 \beta_2\notin N_{K_3/K_2}(K_3^\times).
\tag{A.156.12}
\]
This assertion concerns the specified representative; it does not
exclude another norm preimage of \(\beta_1\).

Here is the full local computation.  At each \(q=11,59\),
the cubic old field \(K_1\) has its degree-one factor and its
unramified quadratic factor.  On the latter the residue of \(\alpha\)
has the indicated equation:
\[
 \begin{array}{c|c|c}
 q&\text{degree-one root of }X^3-3&
          \text{quadratic residue equation for }\alpha\\\hline
 11&9&\alpha^2-2\alpha+4=0\\
 59&49&\alpha^2-10\alpha+41=0.
 \end{array}
\tag{A.156.13}
\]
Both \(K_2/K_1\) and \(K_3/K_2\) are totally ramified cubic
extensions at these branches.  Indeed \(v_{K_1,q}(\rho_1)=8\)
and \(v_{K_2,q}(\rho_2)=8\), prime to three.
The residue fields stay \(\F_{q^2}\), whose multiplicative groups
have three-primary part of order exactly three.

The coefficients \(A_0,A_1\) in (A.156.7) are in
\(\Z[\alpha]\), and \(dA_2\) is in that ring.
Since \(v_{K_2,q}(t)=2\) and \(v_{K_2,q}(q)=3\),
the terms \(A_1t\) and \(A_2t^2\) have positive valuation.
Thus the actual residue of \(\beta_2\) is \(-A_0(\alpha)\).
Substitution of the displayed integer coefficients gives
\[
 \begin{array}{c|c|c|c}
 q&\overline\beta_2&
       \overline\beta_2^{\,(q^2-1)/3}&
       (1-\overline\alpha)^{(q^2-1)/3}\\\hline
 11&6+3\overline\alpha&10+6\overline\alpha&1\\
 59&23+57\overline\alpha&58+6\overline\alpha&1.
 \end{array}
\tag{A.156.14}
\]
The two middle character values are nonidentity cube roots of one;
also \(\overline\beta_2^3=1-\overline\alpha\) in each field,
as required by (A.156.2).
In a totally ramified cubic extension, the residue of the norm
of a unit is its cube.  A norm preimage of the unit \(\beta_2\)
would itself have valuation zero, since the relative residue degree
is one.  The middle column therefore proves that \(\beta_2\)
fails the local norm condition at both of these places, proving
(A.156.12).

For comparison, \(\beta_1=1-\alpha\) meets the unit-norm condition
at these same quadratic places through every higher depth of the
fixed tower.  In a totally ramified extension of degree \(3^j\),
unit norm residues are exactly the \(3^j\)-th powers.
Their image in \(\F_{q^2}^\times\) is the cube subgroup for every
\(j\geq1\), because \(v_3(q^2-1)=1\).
Principal units have all three-power roots in residue characteristic
\(q\ne3\), so a residue norm lifts to an exact unit norm.
The last column of (A.156.14) proves this test for \(\beta_1\).
At the degree-one branches every unit is a three-power scalar norm.
At three use A.140.  At two the distinguished old factor is
\(\Q_2\), with its compatible split root branch.  On the other,
unramified quadratic factor, \(\beta_1\) is a unit and all later
root extensions are unramified, so their unit norms are onto.
At the remaining places use unramified unit norms, and at infinity the
real branch.  Thus the local tests for \(\beta_1\) at all depths
and the failure of this particular \(\beta_2\) are stated on the
same field tower with its actual transition maps.

All residues, cubic characters, coefficient denominators and valuation
bounds in this calculation are independently retained in
\path{code/two_inert_level2_local_norm_certificate.py},
its exact data and its output.
